\documentclass[10pt,a4paper]{amsart}
\usepackage[margin=19mm]{geometry}
\usepackage{amsmath,amssymb,mathtools}
\usepackage[scr=rsfs,cal=euler]{mathalfa}
\usepackage{xfrac}
\usepackage[unicode,hidelinks,bookmarks=false]{hyperref}
\newcommand{\ie}{i.e.\ }
\newcommand{\cf}{cf.\ }
\newcommand{\resp}{respectively\ }
\newcommand{\Subsection}{\subsection}
\providecommand{\nobreakdash}{\nobreak\hbox{-}}
\setlength{\emergencystretch}{2em}
\multlinegap0pt
\usepackage{amssymb}
\usepackage{xcolor}
\newtheorem{corollary}{Corollary}
\newtheorem{theorem}{Theorem}
\title{Cappell-Shaneson knot pairs with the same Alexander polynomial}
\author{Hisaaki Endo, Kazunori Iwaki, Andrei Pajitnov}
\date{Source: arXiv:2604.01045v1 (CC BY 4.0)}
\begin{document}
\maketitle

\noindent\emph{Source and licence.} Cappell-Shaneson knot pairs with the same Alexander polynomial. Version arXiv:2604.01045v1 (CC BY 4.0), distributed by arXiv under the Creative Commons Attribution 4.0 International licence, http://creativecommons.org/licenses/by/4.0/, as recorded in the arXiv licence metadata for this exact version and shown on the abstract page https://arxiv.org/abs/2604.01045v1. Full licence notice: \texttt{licenses/arxiv-2604.01045-CC-BY-4.0.txt}.\par\smallskip

\noindent\textit{Editorial scope.} Counted unit: the article's theorem theorem (source0.tex lines 2398--2443). The article's complete original proof of it is delivered with it (source0.tex lines 2444--2473, one \texttt{proof} environment of 30 lines). No mathematics is written, repaired or re-derived: the statement and the proof are the article's own lines, in the article's own notation, and no retained line is substituted or re-flowed.  Retained uncounted as the source-local context the proof needs: lines 2149--2151.  Nothing was omitted from any retained span, so the package carries no omission record.  No retained line contains a citation, so the wrapper carries no bibliography and no bibliography entry.  No retained line contains a figure environment, an image-inclusion command or drawing code, so no image is delivered and the package declares no asset dependency.  Editorial formatting only, in addition: an AMS \texttt{amsart} preamble that restates the article's own declarations and macros used by the retained lines, namely the environments \texttt{corollary}, \texttt{theorem} on separate counters, together with the standard packages those lines need.  The retained environments print in this document as Corollary 1, Theorem 1 in order of appearance, whereas the article prints the same items with the numbering of its own full document.  The complete native source of the article as distributed in the arXiv e-print for this exact version is bundled unchanged as \texttt{sources/197633/source0.tex}.
    \begin{corollary}\label{corollary:sourceofinffamily}
    Let $p$ be a prime number and $f$ an irreducible monic polynomial with $deg(f) = n$. We denote by $\theta$ a root of $f$. Let us suppose $\bar{f} \in \mathbb{F}_p[x]$ can be factorized into $\Pi_i \bar{g_i}^{e_i}$. Let us suppose $g_k(x)=x-b$. Let $r_k \in \mathbb{Z}[x]$ be a remainder when $f$ is divided by $x - b$ in $\mathbb{Z}[x]$. If $e_k \neq 1$ and $p^2 \nmid r_k$ in $\mathbb{Z}[x]$, then $(p, \theta - b)$ is not invertible and $\{p, \theta - b, \theta ( \theta - b), \ldots , \theta^{n-2} (\theta - b)\}$ is a $\mathbb{Z}$-basis of $(p, \theta -b)$.
    \end{corollary}

\begin{theorem}
    The following two matrices are positive Cappell-Shaneson matrices for any $l \ge 0$. These matrices are not $*$-equivalent and give us inequivalent Cappell-Shaneson knot pairs with the same Alexander polynomial:
    \[
        \begin{bmatrix}
    		0 & 1 & 0 & 0 & 0 & 0 & 0\\
    		0 & 0 & 1 & 0 & 0 & 0 & 0\\
    		0 & 0 & 0 & 1 & 0 & 0 & 0\\
                0 & 0 & 0 & 0 & 1 & 0 & 0\\
                0 & 0 & 0 & 0 & 0 & 1 & 0\\
                0 & 0 & 0 & 0 & 0 & 0 & 1\\
                1 & -G_l & G_l-2 & -G_l^2-2 & G_l^2+1 & -G_l+2 & G_l\\
    		\end{bmatrix},
    \]
    \[
        \begin{bmatrix}
    	-6 & 11 & 0 & 0 & 0 & 0 & 0\\
            0 & 0 & 1 & 0 & 0 & 0 & 0\\
            0 & 0 & 0 & 1 & 0 & 0 & 0\\
            0 & 0 & 0 & 0 & 1 & 0 & 0\\
            0 & 0 & 0 & 0 & 0 & 1 & 0\\
            0 & 0 & 0 & 0 & 0 & 0 & 1\\
            H_l & I_l & J_l & K_l & L_l & M_l & G_l+6\\
    		\end{bmatrix},
    \] with
    \[
    G_l = 121l+20,
    \]
    \[
    H_l = 2012472l^2+1264494l+178211,
    \]
    \[
    I_l = -3689532l^2-2318239l-326720,
    \]
    \[
    J_l = 614922l^2+386353l+54450,
    \]
    \[
    K_l = -102487l^2-64372l-9072,
    \]
    \[
    L_l = 14641l^2+9922l+1445,
    \]
    \[
    M_l = -847l-174.
    \]
\end{theorem}

\begin{proof}
    Define $f_a(x)$ as $x^7-ax^6+(a-2)x^5+(-a^2-1)x^4+(a^2+2)x^3+(-a+2)x^2+ax-1$, which is a positive Cappell-Shaneson polynomial when $a \ge 0$. Let $\theta_a$ be a root of $f_a(x)$. Then, $(11, \theta_{11^2l+20} + 6)$ is not invertible and 
    \[
    \begin{aligned}
    \{\,11,\ &\theta_{11^2l+20} + 6, \ \theta_{11^2l+20} ( \theta_{11^2l+20} + 6), \ \theta_{11^2l+20}^2 (\theta_{11^2l+20} + 6), \\ &\theta_{11^2l+20}^3 (\theta_{11^2l+20} + 6), \ \theta_{11^2l+20}^4 (\theta_{11^2l+20} + 6), \ \theta_{11^2l+20}^5 (\theta_{11^2l+20} + 6)\, \}
    \end{aligned}
    \]
    is a $\mathbb{Z}$-basis of $(11, \theta_{11^2l+20} + 6)$ for any $l \ge 0$ since
    \begin{align*}
        f_{11^2l+20}(x)\equiv &(x+6)^2(x^5+x^4+3x^3+8x+7) \pmod {11} ,\\
        f_{11^2l+20}(x) = &(x+6)(x^6+(-121l-26)x^5+(847l+174)x^4\\
                        &+(-14641l^2-9922l-1445)x^3+(102487l^2+64372l+9072)x^2\\
                        &+(-614922l^2-386353l-54450)x\\
                        &+3689532l^2+2318239l+326720)\\
                        &-22137192l^2-13909434l-1960321
    \end{align*}
    and Corollary~\ref{corollary:sourceofinffamily} can be used. The following matrix corresponds to the ideal class of $(11, \theta_{11^2l+20} + 6)$,
    \[
        \begin{bmatrix}
    	-6 & 11 & 0 & 0 & 0 & 0 & 0\\
            0 & 0 & 1 & 0 & 0 & 0 & 0\\
            0 & 0 & 0 & 1 & 0 & 0 & 0\\
            0 & 0 & 0 & 0 & 1 & 0 & 0\\
            0 & 0 & 0 & 0 & 0 & 1 & 0\\
            0 & 0 & 0 & 0 & 0 & 0 & 1\\
            H_l & I_l & J_l & K_l & L_l & M_l & G_l+6\\
    		\end{bmatrix}.
    \]
    This matrix is not $*$-equivalent to the companion matrix of $f_{11^2l+20}(x)$ since the ideal class corresponding to the companion matrix is invertible. Therefore, the above matrix and the companion matrix give us inequivalent knot pairs that cannot be distinguished by the Alexander polynomial.
\end{proof}

\end{document}
